% ---- pg-485 （书页 473） ----

的任一总体即可.

\noindent\mvbold{证明}\quad 见安德森(文献[2]).\hfill$\blacksquare$

假定所有错分代价都相等(不失一般性, 我们可令所有错分代价为 1), 则最小 ECM 法则就是最小总错分概率法则. 利用导出\hyperref[eq:11-38]{式(11-38)}的类似论据, 我们将 $\mathbf{x}$ 分配到 $\pi_k$, $k=1,2,\cdots,g$, 并使
\phantomsection\label{eq:11-39}
\begin{equation}
\sum_{\substack{i=1\\i\ne k}}^{g}p_if_i(\mathbf{x})
\tag{11-39}
\end{equation}
达到最小. 现在, 当被摒弃的项 $p_kf_k(\mathbf{x})$ 达到最大时, \hyperref[eq:11-39]{式(11-39)}将达到最小. 因此, 当错分代价相等时, 最小 ECM 法则具有下面的简单形式.

\begin{mvframe}[错分代价相同时的最小 ECM 分类法则]
若
\phantomsection\label{eq:11-40}
\begin{equation}
p_kf_k(\mathbf{x})>p_if_i(\mathbf{x}), \text{任意 } i\ne k
\tag{11-40}
\end{equation}
则将 $\mathbf{x}$ 分配到 $\pi_k$, 或等价的, 若
\phantomsection\label{eq:11-41}
\begin{equation}
\ln p_kf_k(\mathbf{x})>\ln p_if_i(\mathbf{x}), \text{任意 } i\ne k
\tag{11-41}
\end{equation}
则将 $\mathbf{x}$ 分配到 $\pi_k$.
\end{mvframe}

值得注意的是, 分类法则(11-40)等价于使“后验”概率 $P(\pi_k\mid \mathbf{x})=P$(给定观测值 $\mathbf{x}$, $\mathbf{x}$ 来自 $\pi_k$)最大化的法则, 其中
\phantomsection\label{eq:11-42}
\begin{equation}
P(\pi_k\mid \mathbf{x})=\frac{p_kf_k(\mathbf{x})}{\sum_{i=1}^{g}p_if_i(\mathbf{x})}=\frac{(\text{先验概率})\times(\text{密度函数})}{\sum(\text{先验概率})\times(\text{密度函数})}\qquad \text{对 } k=1,2,\cdots,g
\tag{11-42}
\end{equation}
\hyperref[eq:11-42]{式(11-42)}是\hyperref[eq:11-9]{式(11-9)}在 $g\ge 2$ 时的推广.

请记住, 一般来说, 最小 ECM 法则有三个要素: 先验概率、错分代价以及密度函数. 这些要素在使用分类法则之前必须指定(或估计出).

\noindent\textbf{例 11.9}\quad (将一个新观测值分到三个已知总体之一)

设虚构的先验概率、错分代价以及密度函数值给定如下表, 让我们将一个观测值 $\mathbf{x}_0$ 分到 $g=3$ 个总体 $\pi_1$, $\pi_2$ 和 $\pi_3$ 中的一个.

\begin{mvtable}
\begin{tabular}{l|cccc}
\toprule
& \multicolumn{4}{c}{真实总体}\\
\cmidrule(lr){2-5}
& & $\pi_1$ & $\pi_2$ & $\pi_3$\\
\midrule
\multirow{3}{*}{将 $\mathbf{x}_0$ 分到:} & $\pi_1$ & $c(1\mid 1)=0$ & $c(1\mid 2)=500$ & $c(1\mid 3)=100$\\
& $\pi_2$ & $c(2\mid 1)=10$ & $c(2\mid 2)=0$ & $c(2\mid 3)=50$\\
& $\pi_3$ & $c(3\mid 1)=50$ & $c(3\mid 2)=200$ & $c(3\mid 3)=0$\\
\midrule
先验概率: & & $p_1=0.05$ & $p_2=0.60$ & $p_3=0.35$\\
密度值: & & $f_1(\mathbf{x}_0)=0.01$ & $f_2(\mathbf{x}_0)=0.85$ & $f_3(\mathbf{x}_0)=2$\\
\bottomrule
\end{tabular}
\end{mvtable}

我们将采用最小 ECM 方法.

$\sum_{i=1}^{3}p_if_i(\mathbf{x}_0)c(k\mid i)$ 的值[见\hyperref[eq:11-38]{式(11-38)}]为
